\documentclass[11pt]{article}
\usepackage[letterpaper, margin=0.85in, top=0.9in, bottom=0.95in, headsep=16pt]{geometry}
\usepackage{libertinus}
\usepackage[scale=0.82]{sourcecodepro}
\usepackage[table]{xcolor}
\usepackage{fvextra, booktabs, longtable, enumitem, titlesec, fancyhdr, microtype, amsmath, array}
\usepackage[most]{tcolorbox}
\usepackage{needspace}
\usepackage[hidelinks]{hyperref}

% J.T.S. palette
\definecolor{ink}{HTML}{1E2B24}
\definecolor{muted}{HTML}{56645B}
\definecolor{accent}{HTML}{2B5E8C}
\definecolor{accentsoft}{HTML}{DCE6EF}
\definecolor{moss}{HTML}{4F6E4C}
\definecolor{mosssoft}{HTML}{E1EADC}
\definecolor{sand}{HTML}{CDBF9E}
\definecolor{surfalt}{HTML}{E7EDE3}
\definecolor{panel}{HTML}{F7F9F5}
\definecolor{rulec}{HTML}{D2DACD}
\definecolor{warn}{HTML}{8A6410}
\color{ink}

\newcommand\Rl[1]{\textcolor{accent}{\textbf{#1}}}
\newcommand\Ms[1]{\textcolor{moss}{\textbf{#1}}}
\newcommand\Dm[1]{\textcolor{muted}{#1}}

\fvset{fontsize=\footnotesize, breaklines, breakanywhere, breakindent=2.5em,
  breaksymbolleft={\tiny\textcolor{sand}{\ensuremath{\hookrightarrow}}}}

\newtcolorbox{termbox}{breakable, enhanced, colback=panel, colframe=rulec, boxrule=0.4pt,
  borderline west={2.2pt}{0pt}{sand}, arc=2pt, left=7pt, right=5pt, top=4pt, bottom=4pt,
  before skip=6pt, after skip=10pt}
\newtcolorbox{thmbox}[1]{breakable, enhanced, colback=accentsoft!45, colframe=accent!55, boxrule=0.5pt,
  arc=2pt, left=8pt, right=8pt, top=5pt, bottom=5pt, before skip=8pt, after skip=10pt,
  title={\sffamily\small\bfseries #1}, coltitle=white, colbacktitle=accent, fonttitle=\sffamily}
\newtcolorbox{keybox}{enhanced, colback=mosssoft!55, colframe=moss!45, boxrule=0.5pt, arc=2pt,
  left=9pt, right=9pt, top=6pt, bottom=6pt, before skip=8pt, after skip=8pt}

\newcommand\cmdtag[1]{\colorbox{accentsoft}{\textcolor{accent}{\ttfamily\bfseries #1}}}
\newcommand\proved{\textcolor{moss}{\textbf{proved}}}
\newcommand\onestep{\textcolor{moss}{\textbf{one step}}}
\newcommand\stepok{\textcolor{accent}{\textbf{step}}}

\titleformat{\section}{\sffamily\Large\bfseries\color{ink}}{}{0pt}{}[\vspace{2pt}{\color{sand}\titlerule[0.8pt]}]
\titlespacing*{\section}{0pt}{20pt}{9pt}
\titleformat{\subsection}{\sffamily\large\bfseries\color{ink}}{}{0pt}{}
\titlespacing*{\subsection}{0pt}{14pt}{5pt}

\pagestyle{fancy}
\fancyhf{}
\renewcommand\headrulewidth{0pt}
\fancyhead[L]{\sffamily\small\color{muted}Quantifier elimination in pvs\_cad}
\fancyhead[R]{\sffamily\small\color{muted}\thepage}

\setlength\parindent{0pt}
\setlength\emergencystretch{3em}
\setlength\parskip{5pt}
\newcommand\code[1]{\mbox{\texttt{#1}}}
% no hyphenation (and no break at a hyphen) inside typewriter identifiers
\makeatletter
\AddToHook{selectfont}{\edef\jts@f{\f@family}\edef\jts@t{\ttdefault}%
  \ifx\jts@f\jts@t\hyphenchar\font=-1\relax\fi}
\makeatother
\renewcommand\arraystretch{1.25}

\begin{document}
\thispagestyle{empty}
\begin{flushleft}
{\sffamily\small\bfseries\color{moss}\MakeUppercase{pvs\_cad \quad·\quad PVS 8.1 \quad·\quad 2 October 2026}}\par\vspace{8pt}
{\fontsize{27}{33}\selectfont\bfseries What quantifier elimination can do}\par\vspace{10pt}
{\large\color{muted} Verified procedures that turn a formula with quantifiers over the reals into an equivalent
formula without them, and a proof command that uses them inside PVS proofs: the theorems, fourteen examples with
the answers as the prover prints them, and what is and is not covered.}
\end{flushleft}
\vspace{4pt}

\begin{keybox}
\textbf{In one line.} Give \texttt{qe8} a formula such as $\exists t.\; t^2 + a t + b = 0$ and it computes a
formula in $a, b$ alone, here a case split on the sign of $b$ that reads
$b<0 \lor b=0 \lor (b>0 \land a^2-4b \ge 0)$, which is equivalent to $a^2 \ge 4b$. That the computed formula has
no quantifiers and holds at exactly the same points as the input is a \textbf{theorem proved in PVS}:
\texttt{qe8\_complete} (and \texttt{qe\_complete} for \texttt{qe}) for every prenex formula with at least one
quantifier and one or more free variables, and \texttt{qelim\_ok} with \texttt{qelim\_qf} for every first-order
formula, with its quantifiers anywhere and any number of free variables, none included.
\end{keybox}

\section{What is proved}

\begin{thmbox}{qe8.qe8\_complete \quad (Tarski--Seidenberg, computed)}
\begin{Verbatim}[fontsize=\small]
qe8_complete: THEOREM cons?(os) AND length(pt) = m IMPLIES
  qf?(qe8(m, os, F, phi)`out) AND
  (qfsem(qe8(m, os, F, phi)`out)(pt) IFF fsem(os, F, phi, pt))
\end{Verbatim}
\end{thmbox}

In words: for every number $m \ge 1$ of free variables (\texttt{m} is a \texttt{posnat}), every prenex formula
\[
  Q_1 x_1 \; Q_2 x_2 \cdots Q_n x_n.\;\; \varphi(y_1,\dots,y_m,\,x_1,\dots,x_n) \qquad (n \ge 1,\; Q_i \in \{\forall,\exists\})
\]
whose matrix $\varphi$ is any Boolean combination of sign conditions ($<0$, $=0$, $>0$, and so $\le,\ge,\ne$) on
polynomials with rational coefficients, and every point $y \in \mathbb{R}^m$: the formula \texttt{qe8} computes is
quantifier-free and holds at $y$ exactly when the input formula does. The hypotheses only fix the shape (at least
one quantifier, \texttt{cons?(os)}; a point with $m$ coordinates): \texttt{qe8} always succeeds and always returns
a quantifier-free formula. \texttt{qe8} runs the Collins CAD (\texttt{docs/collins\_cad.pdf}); \texttt{qe\_all}'s
\texttt{qe\_complete} states the same for \texttt{qe}, the first procedure, which runs the read-closure engine.

\begin{thmbox}{qelim\_ok.qelim\_ok \quad (every first-order formula, quantifiers anywhere)}
\begin{Verbatim}[fontsize=\small]
qelim_qf: LEMMA FORALL phi: FORALL m: qf?(qelim(phi, m))
qelim_ok: THEOREM FORALL phi: FORALL m, pt: length(pt) = m
  IMPLIES (fsem(qelim(phi, m))(pt) IFF fsem(phi)(pt))
\end{Verbatim}
\end{thmbox}
Here \texttt{phi} is any first-order formula over the reals (\texttt{rcf\_fol}'s \texttt{Fm}: \texttt{TRUE}, sign conditions on
polynomials with rational coefficients, \texttt{NOT}, \texttt{AND} and \texttt{EXISTS}, from which the other
connectives and \texttt{FORALL} are built), with its quantifiers anywhere, and $m$ is any natural number, $0$
included. \texttt{qelim} removes the quantifiers from the inside out: under each \texttt{EXISTS} the body is made
quantifier-free first, and then that one quantifier is eliminated by \texttt{qe8}, with the variables in scope as
free variables, or decided by \texttt{decide8} when there are none. No prenex form is needed.
\texttt{gform\_ok}'s \texttt{qe\_g\_ok} states the same for formulas written with PVS's own comparisons and
connectives (\texttt{qe\_g(g, m) = qelim(g2f(g), m)}), and \texttt{fod\_qfd} (every first-order definable set is
quantifier-free definable) follows from \texttt{qelim\_ok}.

{\small
\begin{longtable}{@{}>{\ttfamily}p{0.2\textwidth} p{0.74\textwidth}@{}}
\toprule
\textsf{\small\bfseries Theorem} & \textsf{\small\bfseries What it says}\\
\midrule
qe8\_correct & for every $m\ge1$ and at least one quantifier: \texttt{qe8}'s answer holds at a point of $\mathbb{R}^m$ iff the quantified formula does; \texttt{(cad-qe)} instantiates it for prenex input under \texttt{IMPORTING pvs\_cad}\\
qe8\_isqf & the answer is always quantifier-free (the flag is always set)\\
qe1cc\_bf & one free variable: \texttt{qe1cc}, over one Collins run, always answers, quantifier-free and equivalent (each sector of the free line described exactly, \texttt{epc\_ok})\\
qe2cd\_bf & two or more free variables: over the Collins tower with derivative-closed free levels (\texttt{ctwd}), \texttt{qe2cd} always answers, quantifier-free and equivalent\\
ents\_sep & Thom's lemma at work: when the free families are closed under the derivative, different cells always have different sign patterns\\
qelim\_ok & every first-order formula, quantifiers anywhere, any $m \ge 0$: \texttt{qelim}'s answer, quantifier-free by \texttt{qelim\_qf}, holds exactly where the formula does\\
qe\_g\_ok & the same for formulas written with PVS's comparisons and connectives (\texttt{gform\_def}); \texttt{(cad-qe)} instantiates it for every formula its prenex route does not take (quantifiers inside connectives, binders over subtypes of the reals)\\
qe\_correct,\newline qe\_complete & the statements of \texttt{qe8\_correct} and \texttt{qe8\_complete} for \texttt{qe} (\texttt{qe\_all}), which runs the read-closure engine; \texttt{(cad-qe)} uses \texttt{qe} when the theory imports \texttt{qe\_all} and not \texttt{qe8}\\
qe2c\_ev & \texttt{qe}'s derivative-closed run succeeds for every large enough fuel parameter (the closure stops in a finite universe)\\
\bottomrule
\end{longtable}}

\textbf{What you trust.} The PVS kernel, which checks every step; PVS's ground evaluator, which runs \texttt{qe8},
\texttt{qe\_g} or \texttt{qe} on the concrete input inside the proof (proof by reflection); NASALib's saved proofs
of the theorems the library imports (from \texttt{reals}, \texttt{Sturm}, \texttt{Tarski}, \texttt{structures},
\texttt{analysis}, \texttt{complex}, \texttt{matrices}, \texttt{mult\_poly} and \texttt{interval\_arith}, and
through them from other NASALib libraries such as \texttt{ints}, \texttt{power} and \texttt{trig}), which were not
replayed here; and the reading of the statements: that \texttt{fsem} and \texttt{gsem} (the meaning of
the quantified formula) and \texttt{qfsem} (the meaning of the answer) say what is intended. The procedures and
the strategy code are not trusted; the procedures' correctness is the theorems.

\textbf{Size.} The quantifier elimination described here is 45 theories of the library (\texttt{cad/top.pvs},
2 October 2026) with 550 proved formulas: the first procedure \texttt{qe} (26 theories, with
\texttt{qe\_all} and its examples \texttt{qe\_cad\_ex}), merged answers (\texttt{qe\_mrg\_def}, \texttt{qe\_mrg}),
\texttt{qe8} over the Collins run (13 theories, from \texttt{qe1c\_def} to \texttt{qe8\_ex}), \texttt{qelim}
(\texttt{qelim\_def}, \texttt{qelim\_ok}) and formulas as PVS writes them (\texttt{gform\_def}, \texttt{gform\_ok}).
It is built on the verified CAD library (cylindrical algebraic decomposition, sound and complete decision).

\section{Using it in a proof: \texttt{(cad-qe)}}

\cmdtag{(cad-qe FNUM)} takes a hypothesis or goal with at least one free real term, evaluates the elimination once
inside PVS, prints the answer (\texttt{cad-qe: \dots}), and replaces the formula by the answer. In a theory that
imports \texttt{pvs\_cad}, a prenex formula (one real per quantifier; grouped binders go through the nested form)
is eliminated by \texttt{qe8} and the replacement is justified by \texttt{qe8\_correct}; a formula of any other
shape is eliminated by \texttt{qe\_g}, justified by \texttt{qe\_g\_ok}. In a theory that imports \texttt{qe\_all}
and not \texttt{qe8} (as \texttt{qe\_cad\_ex} does), it evaluates \texttt{qe}, justified by \texttt{qe\_correct}.
Nothing is assumed. A real session (1 October 2026, \texttt{IMPORTING pvs\_cad}; the printed \texttt{cad-qe:} line is
one long line, wrapped here), proving that $t^2+at+b=0$ has a
real root exactly when $a^2\ge 4b$ (the direction ``$\Rightarrow$''):

\begin{termbox}
\begin{Verbatim}[commandchars=\\«»]
w_quad.1 :
{-1}   EXISTS (t: real): t ^ 2 + a * t + b = 0
  |-------
{1}   a ^ 2 >= 4 * b

Rule? \Rl«(cad-qe -1)»
\Ms«cad-qe: ((b < 0) OR ((b = 0) OR ((b > 0) AND ((4*b + (-1)*a^2 = 0) OR»
\Ms«        (4*b + (-1)*a^2 < 0)))))»
length_null rewrites
length[list[mpoly]]((: :))
to 0
Eliminating the quantifiers of -1,
this simplifies to:
w_quad.1 :
{-1}   ((b < 0) OR
        ((b = 0) OR
          ((b > 0) AND
            ((4 * b + (-1) * a ^ 2 = 0) OR (4 * b + (-1) * a ^ 2 < 0)))))
  |-------
{1}   a ^ 2 >= 4 * b
\end{Verbatim}
\end{termbox}
The quantifier is gone; what is left is plain arithmetic, finished with $a^2\ge0$ (3\,s for the
elimination). The other direction is the same command on the goal. The whole lemma is proved in the library:
\texttt{q8\_quad} in \texttt{qe8\_ex} (through \texttt{qe8}) and \texttt{qe\_quad} in \texttt{qe\_cad\_ex} (through
\texttt{qe}).

\section{Examples}
Captured on 1 October 2026 in a theory that imports \texttt{pvs\_cad}; every input here is prenex, so
\texttt{(cad-qe)} eliminates by \texttt{qe8}. Answers are as \texttt{(cad-qe)} prints them, except that parentheses that do
not change the reading (around atoms and around nested disjunctions) are dropped (the transcripts show them). ``Reads as'' is a hand
simplification for the reader; the proved statement is about the printed answer. Times are wall-clock for the
elimination step on an M-series Mac. \proved: the \texttt{IFF} lemma is proved in the library, through
\texttt{qe8} in \texttt{qe8\_ex} (\texttt{q8\_sqrt}, \texttt{q8\_disc}, \texttt{q8\_disc2}, \texttt{q8\_eq2},
\texttt{q8\_quad}) and through \texttt{qe} in \texttt{qe\_cad\_ex} (\texttt{qe\_sqrt}, \texttt{qe\_disc},
\texttt{qe\_disc2}, \texttt{qe\_eq2}, \texttt{qe\_quad}); the first row's lemma, \texttt{qe\_win}, only through
\texttt{qe}. \onestep: the goal is proved by \texttt{(cad-qe)} alone; \stepok: the elimination step is shown (it is
a proved step; closing the arithmetic afterwards was not attempted).

{\small
\begin{longtable}{@{}p{0.255\textwidth} >{\raggedright\arraybackslash}p{0.345\textwidth} p{0.125\textwidth} r l@{}}
\toprule
\textsf{\small\bfseries Input} & \textsf{\small\bfseries Answer, as printed} & \textsf{\small\bfseries Reads as} & \textsf{\small\bfseries Time} & \\
\midrule
\endhead
\multicolumn{5}{@{}l}{\sffamily\small\bfseries\color{moss} One free variable}\\[2pt]
$\exists t\in[0,10].\; (5-t)^2 + (t/2)^2 < D^2$ \newline {\footnotesize\color{muted} a moving point comes within $D$} &
\texttt{125/4 + (-25/4)*D\^{}2 < 0} & $D^2 > 5$ & 3\,s & \proved\\
$\exists t.\; t^2 = D$ & \texttt{(-4)*D = 0 OR (-4)*D < 0} & $D \ge 0$ & 2\,s & \proved\\
$\forall x.\; x^2 + b x + 1 > 0$ & \texttt{4 + (-1)*b\^{}2 > 0} & $b^2 < 4$ & 2\,s & \proved\\
$\exists x\, \exists y.\; x^2 + y^2 < D$ \newline {\footnotesize\color{muted} two bound variables} & \texttt{(-256)*D < 0} & $D > 0$ & 2\,s & \proved\\
$\exists x.\; x > 0 \land x^2 < c$ & \texttt{(-1)*c < 0} & $c > 0$ & 2\,s & \stepok\\
$\forall x\, \exists y.\; y^2 = x^2 + c$ \newline {\footnotesize\color{muted} alternating quantifiers} & \texttt{(-256)*c = 0 OR (-256)*c < 0} & $c \ge 0$ & 2\,s & \stepok\\
$\exists x.\; x^4 - 2x^2 = c$ \newline {\footnotesize\color{muted} degree 4} & \texttt{(-64) + (-64)*c = 0 OR (-64) + (-64)*c < 0} & $c \ge -1$ & 2\,s & \stepok\\
$\exists x.\; x^3 - 3x = c$ & \texttt{TRUE} & always & 2\,s & \onestep\\
\addlinespace
\multicolumn{5}{@{}l}{\sffamily\small\bfseries\color{moss} Two free variables}\\[2pt]
$\exists t.\; t = a \land t = b$ & \texttt{TRUE AND (-1)*b + a = 0} & $a = b$ & 3\,s & \proved\\
$\exists t.\; t^2 + a t + b = 0$ & \texttt{(b < 0) OR (b = 0) OR (b > 0 AND (4*b + (-1)*a\^{}2 = 0 OR 4*b + (-1)*a\^{}2 < 0))} & $a^2 \ge 4b$ & 3\,s & \proved\\
$\exists t.\; t^2 + a t + b < 0$ & \texttt{(b < 0) OR (b = 0 AND 4*b + (-1)*a\^{}2 < 0) OR (b > 0 AND 4*b + (-1)*a\^{}2 < 0)} & $a^2 > 4b$ & 3\,s & \stepok\\
$\forall t.\; t^2 + a t + b \ge 0$ & \texttt{(b = 0 AND 4*b + (-1)*a\^{}2 = 0) OR (b > 0 AND (4*b + (-1)*a\^{}2 > 0 OR 4*b + (-1)*a\^{}2 = 0))} & $a^2 \le 4b$ & 3\,s & \stepok\\
$\exists t.\; t > 0 \land t^2 + a t + b = 0$ \newline {\footnotesize\color{muted} a positive root} & \texttt{(b < 0) OR (b = 0 AND (-2)*a > 0) OR (b > 0 AND ((4*b + (-1)*a\^{}2 < 0 AND (-2)*a > 0) OR (4*b + (-1)*a\^{}2 = 0 AND (-2)*a > 0)))} & $b<0$, or $a<0$ and $a^2\ge4b$ & 4\,s & \stepok\\
\addlinespace
\multicolumn{5}{@{}l}{\sffamily\small\bfseries\color{moss} Three free variables}\\[2pt]
$\exists x.\; a x^2 + b x + c = 0$ \newline {\footnotesize\color{muted} the general quadratic} & 1,187 characters, below & solvable & 8\,s & \stepok\\
\bottomrule
\end{longtable}}

\subsection{A positive root, in the prover}
The same kind of session (the \texttt{cad-qe:} line wrapped):
\begin{termbox}
\begin{Verbatim}[commandchars=\\«»]
s_pos.1 :
{-1}   EXISTS (t: real): t > 0 AND t ^ 2 + a * t + b = 0
  |-------
{1}   b < 0
{2}   (a < 0 AND a ^ 2 >= 4 * b)

Rule? \Rl«(cad-qe -1)»
\Ms«cad-qe: ((b < 0) OR (((b = 0) AND ((-2)*a > 0)) OR ((b > 0) AND»
\Ms«        (((4*b + (-1)*a^2 < 0) AND ((-2)*a > 0)) OR ((4*b + (-1)*a^2 = 0) AND»
\Ms«        ((-2)*a > 0))))))»
length_singleton rewrites
length[list[mpoly]]
      ((: (: mpol((: mpol((: mconst(0), mconst(1) :)) :)),
             mpol((: mconst(0), mconst(1) :)), mconst(1) :) :))
to 1
Eliminating the quantifiers of -1,
this simplifies to:
s_pos.1 :
{-1}   ((b < 0) OR
        (((b = 0) AND ((-2) * a > 0)) OR
          ((b > 0) AND
            (((4 * b + (-1) * a ^ 2 < 0) AND ((-2) * a > 0)) OR
              ((4 * b + (-1) * a ^ 2 = 0) AND ((-2) * a > 0))))))
  |-------
{1}   b < 0
{2}   (a < 0 AND a ^ 2 >= 4 * b)
\end{Verbatim}
\end{termbox}

\subsection{Three free variables: the general quadratic}
$\exists x.\; a x^2 + b x + c = 0$, eliminated in 8\,s. The answer is correct but long: a disjunction of
19 cases over the cells of the decomposition (the sectors of $c$, then of $b$, then of $a$), each described by
sign conditions on $a$, $b$, $c$ and one more polynomial; as a whole it is equivalent to
$(a \ne 0 \land b^2 \ge 4ac) \lor (a = 0 \land (b \ne 0 \lor c = 0))$. In the answer below (1,187 characters,
captured with the table) that polynomial is printed \texttt{((-1)*b\^{}2)*a + (4*c)*a\^{}2}, that is,
$4ca^2 - ab^2 = -a\,(b^2 - 4ac)$: minus $a$ times the discriminant $b^2 - 4ac$.
\begin{termbox}
\VerbatimInput[fontsize=\scriptsize]{s_gen_answer.txt}
\end{termbox}

\section{How it works}
\begin{enumerate}[leftmargin=1.3em, itemsep=3pt]
\item \textbf{One Collins run} on all the variables, free ones outermost (\texttt{qe8}). Its cells form a
cylindrical algebraic decomposition with no certificate (\texttt{col\_cad\_run}); the walk over it, whose
certificate is checked at run time and holds once its effort parameter is large enough (it is raised until then),
gives the truth values at exact sample points. (\texttt{qe} runs the read-closure engine instead.)
\item \textbf{The outermost free variable} is cut into sectors (points and open intervals). Each sector gets an
exact quantifier-free description: a rational endpoint is written directly, an irrational algebraic one by a
polynomial that has it as its only root in a small interval (\texttt{qe\_ep}; proved always possible,
\texttt{epc\_ok}).
\item \textbf{Over each sector}, the cells in the remaining free variables are given the truth value of the
quantified part and the signs of the polynomials there. When equal signs mean equal truth values, the answer for
the sector is the disjunction of the true sign patterns, over a subset of the polynomials that still tells the
cells apart (chosen greedily). With one free variable the sectors themselves are these cells; when their signs do
not tell them apart, the answer is the disjunction of the true sectors' exact descriptions. With one free variable,
either way, the answer is replaced by its merged form, in which each run of consecutive true sectors is one interval,
when that has fewer sign conditions (\texttt{qe\_mrg}); with two or more free variables answers are not merged.
\item \textbf{When they do not} (two or more free variables), \texttt{qe8} runs the Collins tower whose free
families are closed under the formal derivative (\texttt{ctwd}). By Thom's lemma, different cells then always have
different sign patterns (\texttt{ents\_sep}), so the answer is always quantifier-free. (\texttt{qe} instead adds
derivatives to the free families of its run until nothing new appears; that this stops is proved: every
polynomial it can produce lies in one finite set fixed by the input.)
\item \textbf{Formulas of any shape} (\texttt{qelim}, \texttt{qe\_g}): the quantifiers are removed from the inside
out. Under each \texttt{EXISTS} the body is made quantifier-free first; its sign atoms become the family of a
prenex formula with one quantifier, which \texttt{qe8} eliminates (\texttt{decide8} decides it when no variable is
free). \texttt{FORALL} is \texttt{NOT EXISTS NOT}.
\end{enumerate}

\section{Limits}
\begin{itemize}[leftmargin=1.3em, itemsep=3pt]
\item \textbf{Input shape.} First-order formulas over sign conditions on polynomials with rational coefficients.
\texttt{qe} and \texttt{qe8} take prenex formulas with at least one quantifier and one or more free variables
(\texttt{qe\_complete}, \texttt{qe8\_complete}). \texttt{qelim} (\texttt{qelim\_ok}) takes every such formula,
quantifiers anywhere, with any number of free variables, none included; it eliminates from the inside out, so no
prenex conversion is needed. \texttt{(cad-qe)} needs at least one free real term. In a theory that imports
\texttt{pvs\_cad} (or \texttt{gform\_ok}) it takes a formula of any shape through \texttt{qe\_g}
(\texttt{qe\_g\_ok}): quantifiers inside \texttt{NOT}, \texttt{AND}, \texttt{OR},
\texttt{IMPLIES} and \texttt{IFF} (grouped binders included there), and binders over \texttt{real}, \texttt{posreal}, \texttt{nnreal},
\texttt{negreal}, \texttt{npreal} and \texttt{nzreal}; a binder over any other type (\texttt{nat}, \texttt{int},
\texttt{rat}, \texttt{\{x: real | \dots\}}) is refused. A grouped prenex formula goes through its nested form
(\texttt{qe8}, or \texttt{qe} when only \texttt{qe\_all} is imported), with either import; with only \texttt{qe\_all} or \texttt{qe8} imported, \texttt{(cad-qe)} takes
prenex formulas over \texttt{real}.
\item \textbf{Answers are not minimal.} They are disjunctions over cells and sectors; the $a,b,c$ example
shows how long they get. Shortening them further is output polish, not correctness.
\item \textbf{Cost} is the cost of the CAD run: 2 to 8\,s for the examples here; runs with
several bound and free variables together can take minutes (measured with \texttt{qe}: \texttt{QE\_PLAN.md},
section 12).
\item \textbf{Where it is.} In \texttt{cad/}; in the imports of \texttt{cad/top.pvs}, which describes every
theory: \texttt{qe\_ldd} to \texttt{qe\_cad\_ex} (\texttt{qe}, with \texttt{qe\_mrg\_def} and \texttt{qe\_mrg}),
\texttt{qe1c\_def} to \texttt{qe8\_ex} (\texttt{qe8}), \texttt{qelim\_def}, \texttt{qelim\_ok}, \texttt{gform\_def},
\texttt{gform\_ok}.
\end{itemize}

\end{document}
